\vspace{-8mm}
\section{论文发表}

\hypersetup{
    colorlinks=true,
    linkcolor=blue,
    filecolor=magenta,      
    urlcolor=cyan,
    pdftitle={阮炜霖的论文发表},
    pdfpagemode=FullScreen,
}

香港中文大学计算机科学与工程系博士生，研究方向聚焦于多模态与时间序列基础模型、智能体系统、时空学习和具身智能。目前已发表\textbf{7篇CCF-A}和\textbf{2篇CCF-B}级别论文，其中4篇录用论文为第一作者，引用量200+。以下是部分代表性与近期工作：

\begin{enumerate}
\item \textbf{\href{https://openreview.net/forum?id=1WEyvT5Yvg}{ViST: Harnessing Vision Transformation and Reconstruction for Multi-modal Spatio-temporal Forecasting}}
\hfill \textcolor{blue}{第一作者} \\
\textbf{\textcolor{red}{[ACM MM26]}}
\textbf{\textcolor{blue}{(CCF-A)}}
通过视觉变换、多模态条件重建和跨模态融合实现高效时空预测。

\item \textbf{\href{https://arxiv.org/abs/2508.16623}{RAST: A Retrieval Augmented Spatio-Temporal Framework for Traffic Prediction}}
\hfill \textcolor{blue}{第一作者} \\
\textbf{\textcolor{red}{[AAAI26]}} 
\textbf{\textcolor{blue}{(CCF-A)}}
%\textcolor{purple}{\textbf{[ST, RAG]}} 
将检索增强机制与时空建模相结合用于交通预测。

\item \textbf{\href{https://arxiv.org/abs/2602.11550}{TS-Memory: Plug-and-Play Memory for Time Series Foundation Models}}
\hfill \textcolor{blue}{第四作者} \\
\textbf{\textcolor{red}{[KDD26]}}
\textbf{\textcolor{blue}{(CCF-A)}}
通过参数化记忆蒸馏，为冻结的时间序列基础模型提供即插即用、无需在线检索的记忆能力。

\item \textbf{\href{https://arxiv.org/abs/2602.03026}{Visual Reasoning over Time Series via Multi-Agent Systems}}
\hfill \textcolor{blue}{第一作者} \\
\textbf{\textcolor{red}{[Under Review]}} 
%\textcolor{purple}{\textbf{[TS, MAS]}} 
首个面向通用时间序列分析的多智能体系统，实现视觉到数值推理及工具驱动的任务泛化。

\item \textbf{\href{https://arxiv.org/abs/2502.14887}{LDM4TS: Vision-enhanced Time Series Forecasting via Multi-view Latent Diffusion}}
\hfill \textcolor{blue}{第一作者} \\
\textbf{\textcolor{red}{[Under Review]}} 
%\textcolor{purple}{\textbf{[Diffusion, MM]}} 
将时间序列转换为多视图视觉表示，通过频率和文本条件的潜在扩散重建未来模式。

\item \textbf{\href{https://arxiv.org/abs/2508.01727}{OccamVTS: Distilling Vision Models to 1\% Parameters for Time Series Forecasting}}
\hfill \textcolor{blue}{第三作者} \\
\textbf{\textcolor{red}{[AAAI26]}} 
\textbf{\textcolor{blue}{(CCF-A)}}
%\textcolor{purple}{\textbf{[TS, MM]}} 
提出一种知识蒸馏框架，仅从大规模视觉模型中提取1\%的预测时序信息。

\item \textbf{\href{https://proceedings.mlr.press/v267/zhong25a.html}{Time-VLM: Exploring Multimodal Vision-Language Models for Augmented Time Series Forecasting}} \hfill \textcolor{blue}{第二作者} \\
\textbf{\textcolor{red}{[ICML25]}} 
\textbf{\textcolor{blue}{(CCF-A)}}
%\textcolor{purple}{\textbf{[TS, MM, VLM]}} 
利用视觉语言模型通过多模态信息提升时间序列预测。

\item {\small \textbf{\href{https://dl.acm.org/doi/abs/10.1145/3711896.3736962}{Fine-grained Urban Heat Island Effect Forecasting: A Context-aware Thermodynamic Modeling Framework}} \hfill \textcolor{blue}{第二作者} \\
\textbf{\textcolor{red}{[KDD25]}} 
\textbf{\textcolor{blue}{(CCF-A)}}
%\textcolor{purple}{\textbf{[ST, Graph]}} 
用于高分辨率城市热岛预测的上下文感知框架。}

\item \textbf{\href{https://arxiv.org/abs/2404.07919}{STLoRA: Low-rank Adaptation for Spatio-Temporal Forecasting}} \hfill \textcolor{blue}{第一作者} \\
\textbf{\textcolor{red}{[ECML25]}} 
\textbf{\textcolor{blue}{(CCF-B, Oral)}}
%\textcolor{purple}{\textbf{[ST, SFT]}} 
面向时空预测模型的高效低秩自适应技术。
     
\item \textbf{\href{https://ieeexplore.ieee.org/abstract/document/11173261/}{Cross Space and Time: A Spatio-Temporal Unitized Model for Traffic Flow}} \hfill \textcolor{blue}{第一作者} \\
\textbf{\textcolor{red}{[TITS25]}}
\textbf{\textcolor{blue}{(CCF-B, JCR Q1, IF7.9)}}
%\textcolor{purple}{\textbf{[ST]}} 
捕获时空交互的统一方法用于交通预测。

\item \textbf{Towards Multi-Scenario Forecasting of Building Electricity Loads with Multimodal Data} \hfill \textcolor{blue}{第四作者} \\
\textbf{\textcolor{red}{[ACM MM25]}} 
\textbf{\textcolor{blue}{(CCF-A)}}
%\textcolor{purple}{\textbf{[Diffusion]}} 
利用场景条件扩散的跨模态框架用于建筑电力负荷预测。

\item \textbf{\textcolor{red}{[ICCV 2025 Workshop]}} VQualA 2025 Challenge on Face Image Quality Assessment: Methods and Results. \hfill
\item \textbf{\textcolor{red}{[ICLR 2026 Workshop]}} An Empirical Study of the Anchoring Effect in LLMs: Existence, Mechanism, and Potential Mitigations.\hfill
\end{enumerate}
